\chapter[Resultados Esperados]{Resultados Esperados}\label{capitulo6}

Os resultados esperados deste projeto dizem respeito à consolidação de uma solução tecnológica capaz de responder a lacunas identificadas na literatura e na prática psicopedagógica. Na etapa correspondente ao TCC 1, almeja-se a construção de um conjunto consistente de artefatos analíticos e técnicos, suficientes para fundamentar e demonstrar a viabilidade do sistema proposto. Espera-se, portanto, que o trabalho produza uma modelagem conceitual sólida, um protótipo navegável de interfaces e uma especificação de requisitos que traduzam, de forma rigorosa, a lógica da avaliação psicopedagógica aplicada à educação infantil.

Do ponto de vista técnico, espera-se que o sistema modele adequadamente os fluxos de trabalho presentes na avaliação psicopedagógica, contemplando registro de entrevistas iniciais, anamnese, observações, protocolos de instrumentos e devolutivas. A estrutura de dados deve permitir organização longitudinal, associação de registros a metadados e documentação de versões de instrumentos, assegurando rastreabilidade e consistência interna. A expectativa é que a arquitetura proposta ofereça base para implementação escalável, com mecanismos de segurança compatíveis com as exigências da legislação brasileira sobre proteção de dados.

No plano metodológico e profissional, espera-se que a solução concebida contribua para padronização dos registros clínicos, reduzindo fragmentação e inconsistências típicas de anotações exclusivamente manuais. A sistematização proporcionada pelo sistema deve favorecer a continuidade do acompanhamento, permitindo que psicopedagogos acessem, comparem e analisem informações de diferentes momentos da trajetória da criança. De igual modo, espera-se que o modelo de relatórios estruturados contribua para maior clareza nas devolutivas, apoiando tanto as interpretações clínicas quanto a comunicação com famílias e instituições escolares.

Em termos institucionais e éticos, projeta-se que a proposta fomente práticas mais seguras e responsáveis no manejo de dados sensíveis, ao incorporar desde a concepção princípios de minimização, consentimento e controle de acesso. Espera-se que a formalização desses aspectos no projeto sirva de referência para futuras implementações e possa orientar a adoção de procedimentos de governança de dados em contextos clínico-educacionais.

Por fim, antecipa-se que os resultados do TCC 1 estabeleçam bases sólidas para etapas posteriores de desenvolvimento, validação com profissionais e eventual implantação do sistema. A elaboração de artefatos técnicos e conceituais consistentes deverá permitir, em trabalhos futuros, testar empiricamente a ferramenta, avaliar sua utilidade na prática clínica e investigar sua contribuição para análises longitudinais e produção de conhecimento psicopedagógico. A expectativa é que este projeto inicial configure uma proposta tecnicamente fundamentada e metodologicamente coerente, apta a evoluir para uma solução aplicável em contextos reais de avaliação psicopedagógica na educação infantil.
